\section{The two type primes in the denominator}

% BEGIN Y79T1
\paragraph{Y79T1. The local index at a type prime.}
Let $\ell=g_k$ be one of the two type primes and let $i$ be a twisted
localization at $\ell$ on the F78D2 hub. Write $k'$ for the other
index, $R=X.r_i$ and $\rho=X.\rho_i$. Assume the local type-prime
witness for this $i$ when the good-point condition $\mathrm{HG}$ holds:
it gives $\mathrm{GoodR}$ and the pointwise complementary block at
each $\pi\in Z_i$, together with integers $s,S$, signs and jump
sets for $R$, the two required size bounds, and the three conditions
$\mathrm{SignOK},\mathrm{HostJ},\mathrm{SetSQ}$.
The latter conditions use the actual type primes and the original
tame generators. Write $C,C_T$ for the two constants in
$\mathrm{FactsF}$. Then
\[
 \eta_{\rm loc}(i)\le
 \sum_{\pi\in Z_i^{\rm good}}\mathrm{mult}_i(\pi)
                    \,\mathrm{rhs}(C,X.\eta,i,\rho,\pi).
\]

If $\mathrm{HG}$ fails, every good point would yield its witness
by the good-point tie. Hence the good eigenspace sum is zero.
Both defining lattices of $\eta_{\rm loc}$ are sublattices of
this zero space, so the left side and the empty sum both vanish.
If $\mathrm{HG}$ holds, the point-family splitting over the fraction
field gives the requested semisimple decomposition of $R(\pi)$
for each $\pi\in Z_i$: the chosen change of basis identifies
the complement and character blocks, and the diagonal projections
give their summands. This is the point decomposition in DFD6--DFD7
with the semisimplicity from the member's original representation facts.
The size condition $\mathrm{SizeG}$ gives $g_{k'}\le2g_k$, so
$\log g_{k'}\le\log2+\log g_k\le2\log g_k$.
Writing $k_0$ for the source rank parameter, it gives
$g_k\ge(k_0 b)^8$ with $k_0\ge5$ and $b\ge1$.
Thus the type prime is odd. The killing field supplies Selmer
finiteness at every good point and $\mathrm{ShaFinOdd}$ supplies
finiteness for each twist $0<j<t$. These statements and the
specified witness are all the arguments to the paired-character
inequality $\mathrm{aux}_3(C,C_T)$, the second constituent of
$\mathrm{Extra}$ in $\mathrm{FactsF}$. Applying it gives the display.
% END Y79T1

% BEGIN Y79T2
\paragraph{Y79T2. Bounding the terms in the paired-character inequality.}
Suppose $C\ge0$ and let $C_1\ge1$ be the family parameter bound of
A78P3. Put $B=2000(C+1)+5C_1$ and $\eta=L.\mathrm{etaSz}$. Then
for every point in the sum of Y79T1 the right-hand weight is at most
\[
 M_\ell(B)=B x^B\eta+\varepsilon_\ell+
           \sum_{1\le j<t}\mathrm{shaLen}(\ell,j)
            +v_\ell(h_E)+\varepsilon_\ell,
 \qquad \varepsilon_\ell=1/\ell^2 .
\]
Indeed $X.\eta=\lfloor\eta\rfloor\le\eta$ and $1\le t\le\eta$.
Complete splitting at $g_k$ gives $E_w\cong\mathbb Q_\ell$.
On the local cyclotomic extension inversion is an inertia automorphism
on $\mu_\ell$, since this extension is totally ramified. Its image in
the global decomposition group gives cyclotomic value $-1$ modulo
$\ell$. The original ranks have $t\,r$ odd. Thus its $t\,r$-th power
differs from one modulo $\ell$, proving $\nu_L(\ell)=0$ by the
minimum defining $\nu$. The two parts of killing bound the Selmer
length by $\nu_L(\ell)+\varepsilon_\ell=\varepsilon_\ell$ and
the weighted sum at auxiliary primes by $\varepsilon_\ell$.
The other two arithmetic terms in $\mathrm{rhs}$ equal respectively
the displayed sum and class number valuation by their definitions.

For the three remaining terms put $V=n_V=r+t\le2t$, $r\le t$,
$f=2n$ and $v=v_{\rm del}(i)$. The type prime is prime to the host
conductor by its modulus condition, hence outside the host small
set. The original level is one there, so as in C78B1,
$0\le v\le t+v_\ell(t!)\le2t$.
The three terms have sum
\[
 4fV\, C V^2(V+X.\eta+t)+2fr^2+4Vv
       \le (256C+4+16)nt^3\eta .
\]
Use $V+X.\eta+t\le4\eta$ in the first term, and $n,t,\eta\ge1$
in the other two. Now $n,t\le x^{C_1}$, so $nt^3\le x^{4C_1}$.
Since $x\ge1$ and $B\ge\max(256C+20,4C_1)$ the claimed
bound follows. All terms of $M_\ell(B)$ are nonnegative.
% END Y79T2

% BEGIN Y79T3
\paragraph{Y79T3. From the pointwise bound to the original weighted place sum.}
For each allowed lift $j$, each type prime $\ell$ and twisted
coefficient place $w$ of residue characteristic $\ell$, the tie for
$\eta_J$ gives the sum of $\eta_{\rm loc}(i)$ over the localizations
attached to $(j,w)$. These localizations are twisted and have that same
residue characteristic. Apply Y79T1 and Y79T2 to bound the sum by
$M_\ell(B)$ times the sum of their good-point multiplicities.
The $\mathrm{mult}$ field of J79H2 bounds this latter sum by the
local zero-space rank. Its $\mathrm{rank}$ field identifies the rank
with the prescribed $d_{Z,0}(j)$.

Finally the weight of $w$ in the original $\eta_{\rm At}$ is
$\log(\ell)e_w f_w/[L_{\rm coef}:\mathbb Q]$, with $e_w,f_w$
the ramification index and residue degree in the coefficient number
field $L_{\rm coef}$. This follows from $Nw=\ell^{f_w}$ and
the defining degree normalization of the place weight. Each twisted
place contributes once. Summing these nonnegative weights over a
subset of the places above $\ell$ gives at most $\log\ell$, since
$\sum_{w\mid\ell}e_w f_w=[L_{\rm coef}:\mathbb Q]$.
We conclude
\[
 \eta_{\rm At}(L,j,\ell)\le
       (\log\ell)\,d_{Z,0}(j) M_\ell(B).
\]
These are exactly the bounds of $\mathrm{StarH}(L,B)$, including the
two occurrences of $\varepsilon_\ell$. For any $C_s\ge B$,
monotonicity of $u x^u$ for $u\ge0$, $x\ge1$, and
nonnegativity of the other factors give $\mathrm{StarH}(L,C_s)$.
% END Y79T3

% BEGIN Y79T4
\paragraph{Y79T4. Eleven family fields at a common constant and cutoff.}
Use the hypotheses of J79H2, with $C\ge0$ as in the original family
conclusion theorem. Choose $C_1$ by A78P3 before choosing a family
member. In J78H1 choose
\[
 C_s\ge\max\{1,22365,3C_2+1,C_1,C_T,2000(C+1)+5C_1\}.
\]
Use DFD14 with this very $C_s$ when choosing the common cutoff.
It gives $C_s\le b$ there, hence $C_T\le b$.
All other cutoffs of J79H2 can be included in the same maximum.
At every member beyond that cutoff use the original scale guard and
the family premise $\mathrm{TwistedPlaces.ShaFinOdd}(D)$. At this
source revision the latter quantifies over every prime $p$, including
$2$, and requires finite Tate--Shafarevich lengths for $0<e<t$.
Prescribe the arithmetic exceptional set of the member
and take the hub of F78D2. T79W4 supplies the witness of Y79T1 at
both type primes on this hub, using $C_T\le b$; the family size
data supply $\mathrm{SizeG}$. Y79T2 and Y79T3 give
$\mathrm{StarH}(L,C_s)$. The ten fields of J79H2 still hold
because its construction permitted the fixed lower bounds above
before taking the cutoff. Thus on the same $X$ we have
\[
 \begin{gathered}
 \mathrm{pos},\ \mathrm{pass},\ \mathrm{mult},\ \mathrm{rank},\
 \mathrm{kill},\ \mathrm{num},\ \mathrm{bb},\\
 \mathrm{exc},\ \mathrm{typ},\ \mathrm{par},\ \mathrm{small}.
 \end{gathered}
\]
This conclusion uses the paired-character inequality of the original
$\mathrm{FactsF}$. The four outstanding family fields are
$\mathrm{mid},\mathrm{large},\mathrm{two}$ and $\mathrm{ram}$.
% END Y79T4
